% Seminar 2 - Page 3. Fragment for % Seminar 2 - Page 3. Fragment for % Seminar 2 - Page 3. Fragment for % Seminar 2 - Page 3. Fragment for \input{seminars/pages/seminar2-pagina3-en}.
% No preamble and no forced page break.
% EDITORIAL NOTES:
% - The photograph is numbered 3 at bottom right, despite the request
%   saying Page 2. This file therefore represents Page 3.
% - Some quantifier symbols in the image/preimage definitions are overwritten.
% - The example on the right, f:Z->{0,1}, has hard-to-read interval notation.
%   The readings [-1,0] and [1,2] are provisional. Results 1 and 0
%   reflect the manuscript literally, rather than singleton-set notation.
% - Several steps in part (b) are crossed out and have been omitted.
%   Part (a) also implicitly identifies preimages of the same f(x),
%   which need not coincide for a non-injective function.
% - Equality in (b) for an empty indexed family would additionally
%   require f(A)=B; the manuscript does not discuss this exception.
%   Strictly, f(x) in f(intersection) does not force the same x to belong
%   to every X_i without injectivity; a suitable witness is required.
% - The proof ends at the label (c), so no continuation has been supplied.

\subsection*{Seminar 2 - Page 3}

\noindent\fbox{4}\quad Let $f:A\longrightarrow B$, $X\subseteq A$, $Y\subseteq B$.
\[
 f(X)=\{f(x)\mid x\in X\}, \qquad
 f^{-1}(Y)=\{a\in A\mid f(a)\in Y\}.
\]
\noindent (The preimage of $Y$; the fibre of $Y$.)

\medskip
\noindent Let $f:A\longrightarrow B$,
$(X_i)_{i\in I}\subseteq\mathcal{P}(A)$,
$(Y_j)_{j\in J}\subseteq\mathcal{P}(B)$.
\begin{enumerate}
\item[a)]
\[
 f\left(\bigcup_{i\in I}X_i\right)
  =\bigcup_{i\in I}f(X_i).
\]
\item[b)]
\[
 f\left(\bigcap_{i\in I}X_i\right)
  \subseteq\bigcap_{i\in I}f(X_i).
\]
\noindent If $f$ is injective, equality holds.
\item[c)]
\[
 f^{-1}\left(\bigcup_{j\in J}Y_j\right)
  =\bigcup_{j\in J}f^{-1}(Y_j).
\]
\item[d)]
\[
 f^{-1}\left(\bigcap_{j\in J}Y_j\right)
  =\bigcap_{j\in J}f^{-1}(Y_j).
\]
\end{enumerate}

% Example from the right-hand margin (provisional reading).
\noindent\underline{Ex.} $f:\mathbb Z\longrightarrow\{0,1\}$,
\[
 f(x)=\begin{cases}
  0, & x>0,\\
  1, & x\leq 0.
 \end{cases}
 \qquad
 f([-1,0])=1,\quad f([1,2])=0.
\]

\medskip
\noindent\underline{Solution}

\noindent $a)$ Suppose
$f(x)\in f\left(\bigcup_{i\in I}X_i\right)$.
Then $x\in\bigcup_{i\in I}X_i$, so there is some $i_1\in I$
with $x\in X_{i_1}$. Hence
\[
 f(x)\in f(X_{i_1})\subseteq\bigcup_{i\in I}f(X_i).
\]
\noindent This gives the inclusion $\subseteq$.

\smallskip
\noindent Conversely, if
$f(x)\in\bigcup_{i\in I}f(X_i)$, then there is some $i_1\in I$
such that $f(x)\in f(X_{i_1})$. Therefore
\[
 x\in X_{i_1}\subseteq\bigcup_{i\in I}X_i
 \quad\Longrightarrow\quad
 f(x)\in f\left(\bigcup_{i\in I}X_i\right).
\]

\medskip
\noindent $b)$ If
$f(x)\in f\left(\bigcap_{i\in I}X_i\right)$, then
$x\in X_i$ for all $i\in I$, and hence
\[
 f(x)\in f(X_i)\quad(\forall i\in I)
 \quad\Longrightarrow\quad
 f(x)\in\bigcap_{i\in I}f(X_i).
\]
\noindent If $f$ is injective, then
$f(x)\in\bigcap_{i\in I}f(X_i)$ implies
$f(x)\in f(X_i)$ for every $i$, hence
$x\in\bigcap_{i\in I}X_i$.

\medskip
\noindent $c)$
% The proof of (c) resumes on the next manuscript page.
.
% No preamble and no forced page break.
% EDITORIAL NOTES:
% - The photograph is numbered 3 at bottom right, despite the request
%   saying Page 2. This file therefore represents Page 3.
% - Some quantifier symbols in the image/preimage definitions are overwritten.
% - The example on the right, f:Z->{0,1}, has hard-to-read interval notation.
%   The readings [-1,0] and [1,2] are provisional. Results 1 and 0
%   reflect the manuscript literally, rather than singleton-set notation.
% - Several steps in part (b) are crossed out and have been omitted.
%   Part (a) also implicitly identifies preimages of the same f(x),
%   which need not coincide for a non-injective function.
% - Equality in (b) for an empty indexed family would additionally
%   require f(A)=B; the manuscript does not discuss this exception.
%   Strictly, f(x) in f(intersection) does not force the same x to belong
%   to every X_i without injectivity; a suitable witness is required.
% - The proof ends at the label (c), so no continuation has been supplied.

\subsection*{Seminar 2 - Page 3}

\noindent\fbox{4}\quad Let $f:A\longrightarrow B$, $X\subseteq A$, $Y\subseteq B$.
\[
 f(X)=\{f(x)\mid x\in X\}, \qquad
 f^{-1}(Y)=\{a\in A\mid f(a)\in Y\}.
\]
\noindent (The preimage of $Y$; the fibre of $Y$.)

\medskip
\noindent Let $f:A\longrightarrow B$,
$(X_i)_{i\in I}\subseteq\mathcal{P}(A)$,
$(Y_j)_{j\in J}\subseteq\mathcal{P}(B)$.
\begin{enumerate}
\item[a)]
\[
 f\left(\bigcup_{i\in I}X_i\right)
  =\bigcup_{i\in I}f(X_i).
\]
\item[b)]
\[
 f\left(\bigcap_{i\in I}X_i\right)
  \subseteq\bigcap_{i\in I}f(X_i).
\]
\noindent If $f$ is injective, equality holds.
\item[c)]
\[
 f^{-1}\left(\bigcup_{j\in J}Y_j\right)
  =\bigcup_{j\in J}f^{-1}(Y_j).
\]
\item[d)]
\[
 f^{-1}\left(\bigcap_{j\in J}Y_j\right)
  =\bigcap_{j\in J}f^{-1}(Y_j).
\]
\end{enumerate}

% Example from the right-hand margin (provisional reading).
\noindent\underline{Ex.} $f:\mathbb Z\longrightarrow\{0,1\}$,
\[
 f(x)=\begin{cases}
  0, & x>0,\\
  1, & x\leq 0.
 \end{cases}
 \qquad
 f([-1,0])=1,\quad f([1,2])=0.
\]

\medskip
\noindent\underline{Solution}

\noindent $a)$ Suppose
$f(x)\in f\left(\bigcup_{i\in I}X_i\right)$.
Then $x\in\bigcup_{i\in I}X_i$, so there is some $i_1\in I$
with $x\in X_{i_1}$. Hence
\[
 f(x)\in f(X_{i_1})\subseteq\bigcup_{i\in I}f(X_i).
\]
\noindent This gives the inclusion $\subseteq$.

\smallskip
\noindent Conversely, if
$f(x)\in\bigcup_{i\in I}f(X_i)$, then there is some $i_1\in I$
such that $f(x)\in f(X_{i_1})$. Therefore
\[
 x\in X_{i_1}\subseteq\bigcup_{i\in I}X_i
 \quad\Longrightarrow\quad
 f(x)\in f\left(\bigcup_{i\in I}X_i\right).
\]

\medskip
\noindent $b)$ If
$f(x)\in f\left(\bigcap_{i\in I}X_i\right)$, then
$x\in X_i$ for all $i\in I$, and hence
\[
 f(x)\in f(X_i)\quad(\forall i\in I)
 \quad\Longrightarrow\quad
 f(x)\in\bigcap_{i\in I}f(X_i).
\]
\noindent If $f$ is injective, then
$f(x)\in\bigcap_{i\in I}f(X_i)$ implies
$f(x)\in f(X_i)$ for every $i$, hence
$x\in\bigcap_{i\in I}X_i$.

\medskip
\noindent $c)$
% The proof of (c) resumes on the next manuscript page.
.
% No preamble and no forced page break.
% EDITORIAL NOTES:
% - The photograph is numbered 3 at bottom right, despite the request
%   saying Page 2. This file therefore represents Page 3.
% - Some quantifier symbols in the image/preimage definitions are overwritten.
% - The example on the right, f:Z->{0,1}, has hard-to-read interval notation.
%   The readings [-1,0] and [1,2] are provisional. Results 1 and 0
%   reflect the manuscript literally, rather than singleton-set notation.
% - Several steps in part (b) are crossed out and have been omitted.
%   Part (a) also implicitly identifies preimages of the same f(x),
%   which need not coincide for a non-injective function.
% - Equality in (b) for an empty indexed family would additionally
%   require f(A)=B; the manuscript does not discuss this exception.
%   Strictly, f(x) in f(intersection) does not force the same x to belong
%   to every X_i without injectivity; a suitable witness is required.
% - The proof ends at the label (c), so no continuation has been supplied.

\subsection*{Seminar 2 - Page 3}

\noindent\fbox{4}\quad Let $f:A\longrightarrow B$, $X\subseteq A$, $Y\subseteq B$.
\[
 f(X)=\{f(x)\mid x\in X\}, \qquad
 f^{-1}(Y)=\{a\in A\mid f(a)\in Y\}.
\]
\noindent (The preimage of $Y$; the fibre of $Y$.)

\medskip
\noindent Let $f:A\longrightarrow B$,
$(X_i)_{i\in I}\subseteq\mathcal{P}(A)$,
$(Y_j)_{j\in J}\subseteq\mathcal{P}(B)$.
\begin{enumerate}
\item[a)]
\[
 f\left(\bigcup_{i\in I}X_i\right)
  =\bigcup_{i\in I}f(X_i).
\]
\item[b)]
\[
 f\left(\bigcap_{i\in I}X_i\right)
  \subseteq\bigcap_{i\in I}f(X_i).
\]
\noindent If $f$ is injective, equality holds.
\item[c)]
\[
 f^{-1}\left(\bigcup_{j\in J}Y_j\right)
  =\bigcup_{j\in J}f^{-1}(Y_j).
\]
\item[d)]
\[
 f^{-1}\left(\bigcap_{j\in J}Y_j\right)
  =\bigcap_{j\in J}f^{-1}(Y_j).
\]
\end{enumerate}

% Example from the right-hand margin (provisional reading).
\noindent\underline{Ex.} $f:\mathbb Z\longrightarrow\{0,1\}$,
\[
 f(x)=\begin{cases}
  0, & x>0,\\
  1, & x\leq 0.
 \end{cases}
 \qquad
 f([-1,0])=1,\quad f([1,2])=0.
\]

\medskip
\noindent\underline{Solution}

\noindent $a)$ Suppose
$f(x)\in f\left(\bigcup_{i\in I}X_i\right)$.
Then $x\in\bigcup_{i\in I}X_i$, so there is some $i_1\in I$
with $x\in X_{i_1}$. Hence
\[
 f(x)\in f(X_{i_1})\subseteq\bigcup_{i\in I}f(X_i).
\]
\noindent This gives the inclusion $\subseteq$.

\smallskip
\noindent Conversely, if
$f(x)\in\bigcup_{i\in I}f(X_i)$, then there is some $i_1\in I$
such that $f(x)\in f(X_{i_1})$. Therefore
\[
 x\in X_{i_1}\subseteq\bigcup_{i\in I}X_i
 \quad\Longrightarrow\quad
 f(x)\in f\left(\bigcup_{i\in I}X_i\right).
\]

\medskip
\noindent $b)$ If
$f(x)\in f\left(\bigcap_{i\in I}X_i\right)$, then
$x\in X_i$ for all $i\in I$, and hence
\[
 f(x)\in f(X_i)\quad(\forall i\in I)
 \quad\Longrightarrow\quad
 f(x)\in\bigcap_{i\in I}f(X_i).
\]
\noindent If $f$ is injective, then
$f(x)\in\bigcap_{i\in I}f(X_i)$ implies
$f(x)\in f(X_i)$ for every $i$, hence
$x\in\bigcap_{i\in I}X_i$.

\medskip
\noindent $c)$
% The proof of (c) resumes on the next manuscript page.
.
% No preamble and no forced page break.
% EDITORIAL NOTES:
% - The photograph is numbered 3 at bottom right, despite the request
%   saying Page 2. This file therefore represents Page 3.
% - Some quantifier symbols in the image/preimage definitions are overwritten.
% - The example on the right, f:Z->{0,1}, has hard-to-read interval notation.
%   The readings [-1,0] and [1,2] are provisional. Results 1 and 0
%   reflect the manuscript literally, rather than singleton-set notation.
% - Several steps in part (b) are crossed out and have been omitted.
%   Part (a) also implicitly identifies preimages of the same f(x),
%   which need not coincide for a non-injective function.
% - Equality in (b) for an empty indexed family would additionally
%   require f(A)=B; the manuscript does not discuss this exception.
%   Strictly, f(x) in f(intersection) does not force the same x to belong
%   to every X_i without injectivity; a suitable witness is required.
% - The proof ends at the label (c), so no continuation has been supplied.

\subsection*{Seminar 2 - Page 3}

\noindent\fbox{4}\quad Let $f:A\longrightarrow B$, $X\subseteq A$, $Y\subseteq B$.
\[
 f(X)=\{f(x)\mid x\in X\}, \qquad
 f^{-1}(Y)=\{a\in A\mid f(a)\in Y\}.
\]
\noindent (The preimage of $Y$; the fibre of $Y$.)

\medskip
\noindent Let $f:A\longrightarrow B$,
$(X_i)_{i\in I}\subseteq\mathcal{P}(A)$,
$(Y_j)_{j\in J}\subseteq\mathcal{P}(B)$.
\begin{enumerate}
\item[a)]
\[
 f\left(\bigcup_{i\in I}X_i\right)
  =\bigcup_{i\in I}f(X_i).
\]
\item[b)]
\[
 f\left(\bigcap_{i\in I}X_i\right)
  \subseteq\bigcap_{i\in I}f(X_i).
\]
\noindent If $f$ is injective, equality holds.
\item[c)]
\[
 f^{-1}\left(\bigcup_{j\in J}Y_j\right)
  =\bigcup_{j\in J}f^{-1}(Y_j).
\]
\item[d)]
\[
 f^{-1}\left(\bigcap_{j\in J}Y_j\right)
  =\bigcap_{j\in J}f^{-1}(Y_j).
\]
\end{enumerate}

% Example from the right-hand margin (provisional reading).
\noindent\underline{Ex.} $f:\mathbb Z\longrightarrow\{0,1\}$,
\[
 f(x)=\begin{cases}
  0, & x>0,\\
  1, & x\leq 0.
 \end{cases}
 \qquad
 f([-1,0])=1,\quad f([1,2])=0.
\]

\medskip
\noindent\underline{Solution}

\noindent $a)$ Suppose
$f(x)\in f\left(\bigcup_{i\in I}X_i\right)$.
Then $x\in\bigcup_{i\in I}X_i$, so there is some $i_1\in I$
with $x\in X_{i_1}$. Hence
\[
 f(x)\in f(X_{i_1})\subseteq\bigcup_{i\in I}f(X_i).
\]
\noindent This gives the inclusion $\subseteq$.

\smallskip
\noindent Conversely, if
$f(x)\in\bigcup_{i\in I}f(X_i)$, then there is some $i_1\in I$
such that $f(x)\in f(X_{i_1})$. Therefore
\[
 x\in X_{i_1}\subseteq\bigcup_{i\in I}X_i
 \quad\Longrightarrow\quad
 f(x)\in f\left(\bigcup_{i\in I}X_i\right).
\]

\medskip
\noindent $b)$ If
$f(x)\in f\left(\bigcap_{i\in I}X_i\right)$, then
$x\in X_i$ for all $i\in I$, and hence
\[
 f(x)\in f(X_i)\quad(\forall i\in I)
 \quad\Longrightarrow\quad
 f(x)\in\bigcap_{i\in I}f(X_i).
\]
\noindent If $f$ is injective, then
$f(x)\in\bigcap_{i\in I}f(X_i)$ implies
$f(x)\in f(X_i)$ for every $i$, hence
$x\in\bigcap_{i\in I}X_i$.

\medskip
\noindent $c)$
% The proof of (c) resumes on the next manuscript page.
